\subsection{Related work}\label{sec:relatedwork}
The integration of DTs within SoS has become a topic of particular interest, and there is a growing number of secondary studies on the topic.

\textbf{Closest to our work} is the review of \textcite{olsson2023systems} who analyze ten studies in the overlap of DTs and SoS to highlight conceptual challenges, such as integration and interoperability.
Our work provides a systematic treatment of the topic with more breadth ---as we systematically survey the state of the art and eventually analyze 80 primary studies--- and depth ---as we identify recurring architectural strategies, coordination mechanisms, and non-functional properties that characterize emerging SoTS.

Some of the secondary studies related to ours are the following. \textcite{tao2018digital} discuss industrial DT architectures and related patents, which sketch early architectural concepts for coordinating multiple DTs in specific domains such as wind farms. \textcite{verdouw2021digital} introduce six types of DT patterns and illustrates compositions of DTs in agricultural scenarios. Interoperability has been a topic of particular interest recently, with focus on interoperability levels and standards between DTs~\cite{wang2023survey,david2024interoperability}.
Collectively, these works suggest directions toward the composition of complex DT hierarchies but without the element of emergent behavior characteristic to SoS.

Beyond secondary studies, the majority of related primary research focuses on hierarchical integration of DTs to achieve SoS. \textcite{tao2019digital} propose a multi-level DT structure, suggesting that enterprises can achieve SoS by incrementally combining unit-level DTs into complex, higher-order systems. This view is extended by \textcite{gill2024toward}, who advocate for automated horizontal and vertical DT integration and emphasize the need for a unified DT model to enable interoperability across manufacturers. Similarly, \textcite{schroeder2021digital} categorize connectivity levels within DT architectures and demonstrate how aggregated DTs can represent broader system behaviors. \textcite{ghanbarifard2023digital} further reinforce this perspective by discussing distributed DT composition in dynamic, evolving operational contexts. Domain-specific applications of this principle include supply chain integration via Sub-DTs in \textcite{zhang2021supply} and spatial-temporal configurations in \textcite{dietz2020digital}, both modeling SoS-like structures through DT aggregation. 
\textbf{These works highlight the increasing interest in SoTS and motivate empirical inquiries such as ours.}

Despite recent advances, several challenges persist in the realization of SoTS. \textcite{michael2022integration} identify barriers to integrating DTs into SoS, including interoperability gaps, connectivity and privacy issues, and the absence of standardized development practices. These concerns are corroborated by \textcite{semeraro2021digital}, who argue that horizontal data integration is necessary for effective vertical system unification. Further adding to these challenges are the organizational and technical complexities of distributed DTs. \textcite{borth2019digital} outline strategic and architectural difficulties associated with lifecycle coordination, data ownership, and conflicting stakeholder objectives in loosely coupled DT systems. \textbf{These findings highlight the need for unified models and frameworks to advance the state of SoTS, such as our proposal in \secref{sec:classification}.}




%Several systematic literature reviews and systematic mapping studies have been published on digital twins more broadly, synthesizing themes such as DT architectures, lifecycle management, industrial applications, and enabling technologies (\cite{kritzinger2018digital, fuller2020digital, liu2021review, tao2018digital, jones2020characterising, cimino2019review, botin2022digital, barricelli2019survey}). These reviews provide valuable overviews of the state of the art in DT research. However, they primarily consider DTs in isolation or within domain-specific boundaries, such as manufacturing, construction, or healthcare. They do not explicitly address how DTs interact within larger, heterogeneous environments characterized by system of systems principles. Our work targets this gap by systematically analyzing how DTs are composed, coordinated, and scaled within SoS contexts.


RELATED WORK===================================================
\subsection{Related work}\label{sec:relatedwork}

A growing number of secondary studies have investigated digital twins from architectural, technological, and domain-specific perspectives. Broad surveys and systematic reviews synthesize DT concepts, lifecycle models, enabling technologies, and application domains, particularly in manufacturing and industrial settings~\cite{tao2018digital, kritzinger2018digital, fuller2020digital, liu2021review, barricelli2019survey}. While these studies provide valuable overviews of DT research, they predominantly treat DTs as isolated or hierarchically composed systems and do not explicitly analyze their interaction within system-of-systems (SoS) settings.

Closer to our focus are secondary studies that explicitly acknowledge the intersection of DTs and SoS. Olsson et al.~\cite{olsson2023systems} analyze ten studies positioned at this intersection and identify conceptual challenges such as integration and interoperability. However, their review remains exploratory in scope and does not attempt a systematic characterization of architectural strategies, coordination mechanisms, or non-functional concerns across a broad corpus. Similarly, surveys on DT interoperability and standardization~\cite{wang2023survey, david2024interoperability} focus on communication levels, data models, and standards, but stop short of framing DT compositions as SoS or analyzing emergent system behavior.

In contrast to existing secondary studies, our work provides a systematic and structured characterization of systems of twinned systems (SoTS) along multiple dimensions. In terms of \emph{breadth}, we analyze a substantially larger and systematically constructed corpus of primary studies, spanning application domains, architectural styles, coordination approaches, and technological stacks. In terms of \emph{depth}, we go beyond thematic aggregation by identifying recurring architectural strategies and coordination patterns that distinguish SoTS from isolated DTs or purely hierarchical compositions. This enables us to articulate SoTS as a distinct class of systems rather than a loose aggregation of DT integration efforts.

Primary research proposing hierarchical or federated DT architectures~\cite{tao2019digital, schroeder2021digital, gill2024toward, ghanbarifard2023digital} further motivates the need for such a characterization, as these works implicitly adopt SoS principles without systematically analyzing their implications. Our study complements these efforts by providing an empirical synthesis of how such approaches are realized in practice and how they collectively shape the emerging SoTS landscape.

============================================================
Current Related Work

SoS+DT
Secondary/Review/Surveys
olsson2023systems - who analyze ten studies in the overlap of DTs and SoS to highlight conceptual challenges, such as integration and interoperability

Primary Works/Position Papers
gill2024toward - who advocate for automated horizontal and vertical DT integration and emphasize the need for a unified DT model to enable interoperability across manufacturers.

dietz2020digital - spatial-temporal configurations,  modeling SoS-like structures through DT aggregation.

michael2022integration - identify barriers to integrating DTs into SoS, including interoperability gaps, connectivity and privacy issues

borth2019digital - outline strategic and architectural difficulties associated with lifecycle coordination, data ownership, and conflicting stakeholder objectives in loosely coupled DT systems.



DT
Secondary/Review/Surveys
kritzinger2018digital - review of Digital Twins in manufacturing, introducing the Digital Model–Digital Shadow–Digital Twin taxonomy to distinguish levels of physical–digital integration

dalibor2022cross - present a large-scale systematic mapping study of Digital Twins from a software engineering perspective, analysing 356 primary studies to characterise DT counterparts, properties, realization techniques, and evaluation practices, and proposing a feature model to support systematic DT engineering. While the survey explicitly identifies systems-of-systems as a major class of DT counterparts

tao2018digital -  comprehensive state-of-the-art survey of Digital Twin research in industry, framing DTs as cyber–physical systems that integrate modeling, simulation, data fusion, and services across the product lifecycle. While the survey acknowledges interaction and collaboration among multiple virtual models

fuller2020digital - review of Digital Twin enabling technologies, applications, challenges, and open research questions, with a strong emphasis on IoT/IIoT connectivity, data analytics, and Industry 4.0 use cases across manufacturing, healthcare, and smart cities. it frames aggregation primarily as tighter physical–digital coupling and platform integration

negri2017review - review of Digital Twins in Industry 4.0 manufacturing, emphasizing lifecycle continuity, real-time synchronization, and semantic data integration., DTs are modeled at the level of single production systems

semeraro2021digital - present a systematic literature review that synthesizes Digital Twin research across application domains, life-cycle phases, architectures, and enabling technologies. Presents DTs as an emerging cyber physical paradigm. Acknowledges DTs at unit, system, and system-of-systems levels, They argue that horizontal data integration is necessary for effective vertical system unification.

barricelli2019survey - surveys of Digital Twin research, consolidating definitions, core characteristics, application domains, and socio-technical design implications. The survey strongly emphasizes lifecycle continuity, AI-enabled prediction, and human–DT interaction

ferko2023standardisation - systematically assess how existing Digital Twin architectures in manufacturing align with the ISO 23247 reference architecture, showing that most implementations realize only a narrow subset of the prescribed functional entities and largely neglect cross-system, interoperability, and lifecycle concerns.

ferko2022architecting - present a systematic mapping study of software architectures for Digital Twins, analysing 140 primary studies to characterise architectural solution types, quality attributes, and recurring architectural patterns.


% 
% xu2023survey - survey of Digital Twin–enabled Industrial IoT, framing DTs as high-fidelity cyber-physical replicas that integrate control, communication, and computing (3C), and that are increasingly augmented with AI and blockchain for intelligent and secure operation. largely assumes unified system ownership and control, and does not address the composition of multiple autonomous digital twins

% wang2023survey - survey IoDT as a large-scale, distributed extension of cyber–physical systems, emphasizing inter-twin cooperation, decentralized data exchange, and security governance. However, despite this decentralization rhetoric, DT interactions are largely assumed to operate within a single system boundary, stopping short of treating DTs as independently governed constituents in a true system-of-systems sense.

% ghanbarifard2023digital - discussing distributed DT composition in dynamic, evolving operational contexts, surveying use cases, challenges, and enabling technologies across large-scale, highly interconnected systems such as maritime operations, smart cities, and industrial infrastructures.



Primary Works/Position Papers
verdouw2021digital - introduce six types of DT patterns and illustrates compositions of DTs in agricultural scenarios

tao2019digital - analyzes CPS and DTs from multiple perspectives. propose a multi-level DT structure, suggesting that enterprises can achieve SoS by incrementally combining unit-level DTs into complex, higher-order systems.

schroeder2021digital - categorize connectivity levels within DT architectures and demonstrate how aggregated DTs can represent broader system behaviors. 



SoS -------------------------
Secondary/Review/Surveys
gorod2008system - review the SoS literature to illustrate the need to create an SoSE management framework based on the demands of constant technological progress in a complex dynamic environment. Discuses composing complex systems in a general sense. 

uday2015designing - present a comprehensive survey of metrics, methods, and design challenges for engineering resilience in Systems-of-Systems, critically examining the limitations of traditional reliability and risk approaches and synthesizing multidisciplinary strategies for coping with disruption, recovery, and uncertainty. While the survey draws heavily on cyber-physical infrastructure domains and networked operational systems, it treats resilience at the architectural and system-interaction level and does not explicitly frame these mechanisms in terms of Digital Twins or continuous physical–digital synchronization.

nielsen2015systems - present a comprehensive survey of System-of-Systems engineering that synthesizes foundational concepts, categorizes SoS types, and systematically maps eight core SoS dimensions. Discusses cyber physical example domains for SoS 

keating2003system - present an early foundational treatment of System-of-Systems Engineering, explicitly framing SoSE as the engineering of integrated complex metasystems composed of autonomous, heterogeneous complex systems operating under uncertainty, emergence, and evolution. While the paper discusses large-scale sociotechnical and infrastructure systems that combine physical assets, computation, and human decision-making, it predates and does not explicitly adopt the cyber-physical systems paradigm or Digital Twin concepts.

klein2013systematic - do a lit review on systems of system architecture, finding most solutions are fragmented and there no strong convergence towards specific architecture or tooling. Focuses specifically on the architecture dimension of SoS — how SoS structures are described and classified in literature. Helps identify architectural patterns, gaps, and trends across published work.

axelsson2015systematic - systematic mapping of System-of-Systems Engineering publications, characterizing research topics, application domains, and community structure. the study explicitly references cyber-physical systems and IoT as closely related and future application domains and highlights architecture, modeling, and integration as dominant themes,


% 
% fang2021system - review on the SoS architecture selection to preserve the dynamic properties of SoS. Mentions The recent digital twin technologies might offer a great opportunity to collect adequate high-quality architecture data that could be used to be analyzed with different methods for addressing different issues during SoS architecting,  in their discussion on future opportunities. Along with an empahsis on model driven and ML incoportations.


% santos2022evaluation - present a systematic review of methods for evaluating System-of-Systems software architectures, analyzing modeling techniques, quality attributes, and the maturity of existing evaluation approaches for dynamic and evolving SoS. the survey discusses cyber-physical application domains and runtime architectural reconfiguration

% dridi2020system - literature review of System-of-Systems modelling approaches, surveying challenges, quality attributes, and a wide range of MBSE-based techniques (e.g., MDA, ADLs, SOA, ontologies) used to describe static and dynamic SoS architectures. While the review frequently references cyber-physical application domains and emphasizes architectural evolution and reconfiguration, it remains focused on design-time modelling abstractions

% vargas2016approaches - systematic review of System-of-Systems integration research, focusing on software engineering methods, interoperability mechanisms, and architectural approaches used to integrate heterogeneous constituent systems. While the review emphasizes integration challenges in distributed and often cyber-physical application domains and references service-oriented and event-based architectures





% ==================================================================================
\subsection{Related work}\label{sec:relatedwork}
This section reviews prior work relevant to this study in SoS engineering, DT research, and recent efforts that combine these perspectives. 

\paragraph{System of Systems Engineering}
SoS engineering research establishes the conceptual foundations for composing autonomous, heterogeneous systems into larger capabilities. Foundational work by \citet{keating2003system} frames SoS as independently governed systems operating under uncertainty, emergence, and evolution. Subsequent surveys by \citet{gorod2008system} and \citet{nielsen2015systems} consolidate these ideas by categorizing SoS types and dimensions and by identifying recurring coordination and interoperability challenges. In doing so, these studies frequently reference cyber-physical application domains, including networked infrastructure, transportation, and defense systems.
More focused secondary studies examine specific engineering concerns within these domains. For example, \citet{uday2015designing} discuss disruption and recovery in system of systems and reference cyber-physical infrastructures, such as urban and critical infrastructure systems, as representative application contexts. Architectural reviews and mappings by \citet{klein2013systematic} and \citet{axelsson2015systematic} investigate SoS structures, modeling approaches, and integration strategies, while highlighting fragmentation and limited convergence in existing architectural solutions. Related reviews of SoS modeling and architectural evaluation similarly draw on cyber-physical use cases and dynamic reconfiguration scenarios \citep{dridi2020system, santos2022evaluation, vargas2016approaches}, but remain centered on interaction-level coordination.
Overall, SoS research uses cyber-physical systems as motivating contexts but does not provide concrete mechanisms for representing, synchronizing, or managing the evolving state of constituent systems during operation.


\paragraph{Digital Twin Research}
DT research has been extensively reviewed across manufacturing and industrial contexts. Surveys by \citet{kritzinger2018digital} and \citet{negri2017review} characterize DTs as cyber--physical systems that emphasize lifecycle continuity, real-time data synchronization, and semantic integration, while typically modeling DTs at the level of individual machines or production systems. Broader surveys by \citet{tao2018digital} and \citet{fuller2020digital} further reinforce this cyber--physical perspective by framing DTs as integrating modeling, simulation, and analytics across product lifecycles and application domains.
Several large-scale secondary studies move beyond single-system perspectives and explicitly consider more complex DT counterparts. A systematic mapping by \citet{dalibor2022cross} analyzes DTs from a software engineering perspective and identifies SoS as a major class of DT counterparts. Similarly, \citet{semeraro2021digital} synthesize DT research across architectures and lifecycle phases and acknowledge DTs at unit, system, and SoS levels, arguing that horizontal data integration is necessary for effective system unification. Related studies propose multi-level or aggregated DT structures \citep{tao2019digital, verdouw2021digital, schroeder2021digital}, illustrating how unit-level DTs may be combined to represent broader system behavior.
Complementary architectural and standardization-focused work by \citet{ferko2023standardisation} and \citet{ferko2022architecting} makes this limitation explicit by showing that most DT implementations realize only a subset of the functional entities outlined by standards such as ISO~23247. These studies emphasize the need for stronger interoperability, cross-system integration, and lifecycle coordination, which remain insufficiently addressed in current DT architectures.
Overall, Digital Twin research has increasingly focused on composition and aggregation to enable cross-system interaction and more complex structures, motivating closer examination of how such compositions are architected and coordinated.



\textbf{Closest to our work is the review of} \citet{olsson2023systems}, which analyzes ten studies at the DT--SoS intersection and identifies recurring conceptual challenges related to integration, interoperability, and coordination across multiple Digital Twins. Their work highlights the lack of shared conceptual models for aligning DT-centric and SoS-centric perspectives. Our study builds on this foundation by complementing conceptual synthesis with an analysis of how SoS properties are reflected in DT-based architectures and how these properties influence coordination and interoperability in practice. In particular, we provide a systematic treatment of the topic that combines breadth, through a comprehensive survey of the state of the art, and depth, by identifying recurring architectural strategies, coordination mechanisms, and non-functional properties that characterize emerging systems of twinned systems.

Several other studies address specific aspects of DT-based SoS. \citet{dietz2020digital} investigate spatial--temporal DT aggregation, modeling SoS-like structures through the composition of multiple Digital Twins. \citet{gill2024toward} advocate for automated horizontal and vertical DT integration and emphasize unified DT models to enable interoperability across manufacturers and organizational boundaries. Other studies focus on barriers and complexities in DT-based SoS. \citet{michael2022integration} identify interoperability gaps, connectivity and privacy issues, and the lack of standardized development practices when integrating DTs into SoS. Complementing this perspective, \citet{borth2019digital} highlight strategic and architectural challenges related to lifecycle coordination, data ownership, and conflicting stakeholder objectives in loosely coupled DT ecosystems.





